\subsection{G 的闭子群上的 Haar 测度所成的空间}

引理 1. --- 设 \(\alpha\) 为 G 上一个 \(\neq 0\) 的正测度，S 为它的支集；下列两个条件等价：

\begin{enumerate}
\def\labelenumi{\alph{enumi})}
\item
  S 是 G 的一个闭子群，且 \(\alpha\) 在 S 上诱导的测度是 S 上的一个右 Haar 测度。
\item
  \(\delta(s)\alpha = \alpha\) 对每一 \(s \in S\) 成立。
\end{enumerate}

此外，当这些条件满足时，\(t \in \mathbf{G}\) 中使得 \(\delta(t)\alpha = \alpha\) 的那些 t 组成的集合等于 \(S\)。

显然 a) 蕴含 b)；反之，关系 b) 蕴含 \(Sx = S\)（对每一 \(x \in S\)）；换句话说，关系 \(x \in S\) 与 \(y \in S\) 蕴含 \(y \in Sx\)，或者再说一次 \(yx^{-1} \in S\)，而由于 \(S\) 非空，\(S\) 是 \(G\) 的一个闭子群。于是 \(t \in G\) 中使得 \(St = S\) 的那些 t 组成的集合等于 \(S\) 自身，由此得出最后的断言。

在本节其余部分，我们用 \(\Gamma\) 表示 G 上满足引理 1 的条件的 \(\neq0\) 正测度组成的集合，且对每一 \(\alpha\in\Gamma\)，用 \(H_{\alpha}\) 表示作为 \(\alpha\) 的支集的 G 的闭子群。

命题 1. --- 集合 \(\Gamma\) 在赋有淡拓扑的空间 \(\mathcal{M}_{+}(G) = \{0\}\) 中是闭的。

我们先证明下列引理：

引理 2. --- 设 X 为一个局部紧空间，且对每一测度 \(\alpha \in \mathcal{M}_{+}(X) - \{0\}\)，设 \(S_{\alpha}\) 为 \(\alpha\) 的支集。设 \(\Phi\) 为 \(\mathcal{M}_{+}(X) - \{0\}\) 上的一个滤子，它淡收敛于一个测度 \(\alpha_{0} \neq 0\)。则对 \(\alpha_{0}\) 的支集的一点 s 的每个邻域 V，存在一个集合 \(M \in \Phi\)，使得对每一 \(\alpha \in M\)，有 \(V \cap S_{\alpha} \neq \varnothing\)。

因为，若 \(\varphi \in \mathcal{K}_{+}(\mathbf{X})\) 是一个支集包含在 \(\mathbf{V}\) 中且满足 \(\int \varphi(x)d\alpha_0(x) > 0\) 的函数，则按定义存在一个集合 \(\mathbf{M} \in \Phi\)，使得 \(\int \varphi(x)d\alpha(x) > 0\)（对所有 \(\alpha \in \mathbf{M}\)），这就蕴含 \(\mathbf{V} \cap \mathbf{S}_{\alpha} \neq \emptyset\)。

引理 3. --- 设 E 为一个由滤子 \(\Phi\) 加滤的集合，且设 \(\xi \mapsto \alpha(\xi)\) 为 E 到 \(\Gamma\) 中的一个映射，它关于 \(\Phi\) 淡收敛于一个测度 \(\alpha_0 \neq 0\)。另一方面，设 \(\xi \mapsto t_{\xi}\) 为 E 到 G 中的一个映射，使得 \(t_{\xi} \in \mathrm{H}_{\alpha(\xi)}\)（对每一 \(\xi \in \mathrm{E}\)）。若 \(s\) 是映射 \(\xi \mapsto t_{\xi}\) 关于 \(\Phi\) 的一个聚点，则 \(\delta(s)\alpha_0 = \alpha_0\)。

若有必要，把 \(\Phi\) 换成一个更细的滤子，我们可以假设 \(s\) 是 \(\xi \mapsto t_{\xi}\) 关于 \(\Phi\) 的一个极限；由引理 1，\(\delta(t_{\xi})\alpha(\xi) = \alpha(\xi)\)（对每一 \(\xi \in E\)），而结论由映射 \((u,\lambda) \mapsto \delta(u)\lambda\) 在 \(G \times \mathcal{M}_{+}(G)\) 上的连续性得出（§3, No.~3 命题 13）。

为证明命题 1，由引理 1，只需证明：若一个滤子 \(\Psi\)（\(\Gamma\) 上的）淡收敛于一个测度 \(\alpha_0 \neq 0\)，且 \(s\) 属于 \(\alpha_0\) 的支集，则 \(\delta(s)\alpha_0 = \alpha_0\)。现在，对每一邻域 \(V\)（\(s\) 在 \(G\) 中的），由引理 2，存在一个 \(M \in \Psi\)，使得对每一 \(\alpha \in M\)，有 \(V \cap H_\alpha \neq \varnothing\)。对每一邻域 \(V\)（\(s\) 的）与每一 \(\alpha \in \Gamma\)，设 \(t_{V,\alpha}\) 为 \(V \cap H_\alpha\) 的一点（若 \(V \cap H_\alpha \neq \varnothing\)），在相反情形设为 \(H_\alpha\) 的任意一点；若 \(\Theta\) 是 \(s\) 的邻域滤子的截滤子，且 \(\Phi\) 是积滤子 \(\Theta \times \Psi\)，则由上所述，\(s\) 是 \((V,\alpha) \mapsto t_{V,\alpha}\) 关于 \(\Phi\) 的一个聚点。由于另一方面映射 \((V,\alpha) \mapsto \alpha\) 以 \(\alpha_0\) 为极限（关于 \(\Phi\)），命题由引理 3 得出。

命题 2. --- 设 \(\varphi\) 为 \(\mathcal{K}_{+}(\mathbf{G})\) 中的一个函数，使得 \(\varphi(e) > 0\)。则集合 \(\Gamma_{\varphi}\)（即测度 \(\alpha \in \Gamma\) 中满足 \(\int \varphi(x) d\alpha(x) = 1\) 的那些）对淡拓扑是紧的。

集合 \(\Gamma_{\varphi}\) 是 \(\Gamma\) 与 \(\mathcal{M}(\mathrm{G})\) 中超平面（由 \(\alpha\)——满足 \(\int \varphi(x)d\alpha(x)=1\)——形成）的交；由于这个超平面在 \(\mathcal{M}(\mathrm{G})\) 中是淡闭的且不含 0，由命题 1 可知 \(\Gamma_{\varphi}\) 在 \(\mathcal{M}(\mathrm{G})\) 中是淡闭的。因此只需证明：对 G 的每一紧子集 K，有 \(\sup_{\alpha\in\Gamma_{\varphi}}\alpha(\mathrm{K})<+\infty\)（Ch. III, §1,

No.~9, 命题 15)。现在，设 U 为 e 在 G 中由不等式 \(\varphi(x) > \varphi(e)/2\) 定义的开邻域；由于 \(1 = \int \varphi(x) d\alpha(x) \geqslant \int_{\mathrm{U}} \varphi(x) d\alpha(x)\)（对 \(\alpha \in \Gamma_{\varphi}\)），可见在置 \(c = 2/\varphi(e)\) 后，有 \(\alpha(\mathrm{U}) \leqslant c\)（对每一 \(\alpha \in \Gamma_{\varphi}\)）。设 V 为 e 在 G 中的一个对称开邻域，使得 \(V^{2} \subset U\)；我们来证明 \(\alpha(Vx) \leqslant c\)（对每一 \(x \in G\) 与每一 \(\alpha \in \Gamma_{\varphi}\)）。事实上，若 Vx 不与支集 \(H_{\alpha}\)（\(\alpha\) 的）相交，这一关系是平凡的；反之，若存在一个 \(h \in Vx \cap H_{\alpha}\)，则对某个 \(v \in V\) 有 h = vx，从而

\[
\mathbf {V} x = \mathbf {V} v ^ {- 1} h \subset \mathbf {V} ^ {2} h \subset \mathrm{U} h,
\]

而由于 \(\delta(h)\alpha = \alpha\)，得出 \(\alpha(\mathrm{V}x) \leqslant \alpha(\mathrm{U}h) = \alpha(\mathrm{U}) \leqslant c\)。现在设 \((x_i)_{1 \leqslant i \leqslant n}\) 为 K 的点的一个序列，使得诸 \(Vx_i\) 构成 K 的一个覆盖；由上所述，\(\alpha(K) \leqslant \sum_{i=1}^{n} \alpha(Vx_i) \leqslant nc\)（对每一 \(\alpha \in \Gamma_\varphi\)）；证毕。

命题 3. --- 在命题 2 的假设下，映射 \(\alpha \mapsto \left( \langle \varphi, \alpha \rangle, \frac{\alpha}{\langle \varphi, \alpha \rangle} \right)\) 是 \(\Gamma\) 到积空间 \(\mathbf{R}_{+}^{*} \times \Gamma_{\varphi}\) 上的一个同胚。

由于映射 \(\alpha \mapsto \langle \varphi, \alpha \rangle\) 是淡连续的，只需注意到 \(\langle \varphi, \alpha \rangle \neq 0\) 对每一测度 \(\alpha \in \Gamma\) 成立，因为 \(e\) 属于支集 \(\mathbf{H}_{\alpha}\)（\(\alpha\) 的）且 \(\varphi(e) > 0\)。

